
\needspace{3\baselineskip}
\color{uniblue}{\subsection{Коммутационные аппараты (КА)}}
\color{black}

\begin{enumerate}[label=\arabic{section}.\arabic{subsection}.\arabic{enumi}, labelsep=4pt, leftmargin=0pt, itemindent=57pt, itemsep=0pt, parsep=5pt] % если что можно крутить itemsep

%%%%%%%%%%%%%%%%%%%%%%%%%%%%%%%%%%%%%%%%%%%%%%%%%%%%%%%%%%%%%%%%%%%%% ОБЩИЕ СВЕДЕНИЯ %%%%%%%%%%%%%%%%%%%%%%%%%%%%%%%%%%%%%%%%%%%%%%%%%%%%%%%%%%%%%%

\item Общие сведения

\begin{enumerate}[label=\arabic{section}.\arabic{subsection}.\arabic{enumi}.\arabic*, labelsep=4pt, leftmargin=0em, itemindent=65pt, parsep=0pt]

\item
Функциональный блок <<КА>> предназначен для реализации функций контроля и управления высоковольтными коммутационными аппаратами (заземляющий нож, разъединитель) в составе защищаемого присоединения.  
\item
Функциональный блок <<КА>> предусматривает контроль и управление следующими коммутационными аппаратами:
\begin{itemize}
    \item заземляющий нож;
    \item выкатной элемент.
\end{itemize}

\item
Ввод функционального блока <<КА>> в работу на этапе параметрирования устройства осуществляется программным переключателем SGF1 <<Ввод\_функции>>. Введенное состояние функционального блока соответствует активному состоянию сигнала <<КА: Ввод>>.
\item
Функциональный блок <<КА>> может быть выведен из работы оперативно путем активации сигнала <<ОВ КА>>, о чем сигнализирует активное состояние сигнала <<КА: Оперативный вывод>>.
\item
Логика управления режимами работы функционального блока <<КА>> выполнена в соответствии с рисунком \ref{ka:img_gen}. В таблице \ref{ka:tbl_gen} приведены параметры, необходимые для настройки функционального блока <<КА>>. 

\vspace{2mm}
\begin{figure}[H]
\centering
\includegraphics[width=1\textwidth,height=1\textheight,keepaspectratio]{\fbpath/213. КА ТН/_latex/img/1.pdf}
\captionof{figure}{Логика управления режимами функционального блока <<КА>>}\label{ka:img_gen}
\end{figure}

\small
\begin{longtable}{|>{\centering\arraybackslash}m{5.3cm}|>{\centering\arraybackslash}m{3.3cm}|>{\centering\arraybackslash}m{4.2cm}|>{\centering\arraybackslash}m{1.8cm}|>{\centering\arraybackslash}m{1cm}|}
\caption{Параметры для настройки функционального блока <<КА>>\hfill\vspace{-0.5\baselineskip}}\label{ka:tbl_gen}\\ 
\hline
\rowcolor{gray!30}
Параметр (Параметр на ИЧМ) & Условное обозначение на схеме & Значение/ Диапазон & Единица измерения & Шаг \\ 
\hline
\endfirsthead
\caption*{\hspace{3pt}\emph{Продолжение таблицы \ref{ka:tbl_gen}\hfill\vspace{-0.5\baselineskip}}} \\ % сделано по ГОСТ 2.105 п.6.8.7
\hline
\rowcolor{gray!30}
Параметр (Параметр на ИЧМ) & Условное обозначение на схеме & Значение/ Диапазон & Единица измерения & Шаг \\ 
%\hline
\endhead
%\hline
\endfoot
%\hline
\endlastfoot
%==+t1*LLN0|VT_SwitchDevice
\centering Ввод функции в работу (Ввод\_функции) & \centering SGF1 & \centering 0 = Не предусмотрено\\1 = Предусмотрено & \centering -- & \centering \arraybackslash -- \\
\hline
%===t1
\end{longtable}
\normalsize

\end{enumerate}

%%%%%%%%%%%%%%%%%%%%%%%%%%%%%%%%%%%%%%%%%%%%%%%%%%%%%%%%%%% ЗН %%%%%%%%%%%%%%%%%%%%%%%%%%%%%%%%%%%%%%%%%%%%%%%%%%%%%%%%%%

\item Заземляющий нож (ЗН)

\begin{enumerate}[label=\arabic{section}.\arabic{subsection}.\arabic{enumi}.\arabic*, labelsep=4pt, leftmargin=0em, itemindent=65pt, parsep=0pt]

\item
При наличии входного сигнала <<УЗН: Отключить>> будет формироваться команда на отключение заземляющего ножа <<ЗН: Отключить (реле)>> на внутреннее выходное реле устройства с целью коммутации выходной цепи заземляющего ножа. С помощью уставки SGF2 <<Режим\_откл>> задается режим отключения заземляющего ножа (длительный/импульсный). В импульсном режиме выдаваемый импульс длится на протяжении времени, задаваемого уставкой Т2. При наличии отключенного состояния заземляющего ножа <<ЗН: Отключен>> произойдет сброс команды отключения как в импульсом, так и в длительном режимах. Сброс команды отключения также осуществляется с помощью команды <<Сброс>>.
\item
При наличии входного сигнала <<УЗН: Включить>> будет формироваться команда на включение заземляющего ножа <<ЗН: Включить (реле)>> на внутреннее выходное реле устройства с целью коммутации выходной цепи заземляющего ножа. С помощью уставки SGF3 <<Режим\_вкл>> задается режим включения заземляющего ножа (длительный/импульсный). В импульсном режиме, выдаваемый импульс, длится на протяжении времени, задаваемого уставкой Т3. При наличии включенного состояния заземляющего ножа <<ЗН: Включен>> произойдет сброс команды включения как в импульсом, так и в длительном режимах. Сброс команды включения также осуществляется с помощью команды <<Сброс>>.
\item
При выводе функции из работы (отсутствие выходного сигнала <<ЗН: Ввод>>) происходит сброс команд управления заземляющим разъединителем.
\item
В логике функции реализовано формирование выходных сигналов от тележки ТН без выдержки времени таких как:
\begin{itemize}
    \item промежуточное положение заземляющего разъединителя <<ЗН: Промеж. положение>> (формируется если одновременно присутствуют/отсутствуют оба сигнала <<ЗН отключен (б/к)>>, <<ЗН включен (б/к))>>;
    \item отключенное положение заземляющего разъединителя <<ЗН: Отключен>>;
    \item включенное положение заземляющего разъединителя <<ЗН: Включен>>.
\end{itemize}
\item
Неисправное положение заземляющего разъединителя <<ЗН: Неиспр. положение>> формируется по истечении регулируемой выдержки времени задаваемой уставкой Т1.

\item
Ввод функции <<ЗН>> в работу на этапе параметрирования устройства осуществляется программным переключателем SGF1 <<Ввод\_функции>>. Введенное состояние функции соответствует активному состоянию сигнала <<ЗН: Ввод>>.

\item
Алгоритм работы функции <<ЗН>> выполнен в соответствии с рисунком \ref{ka:img_zn}. В таблице \ref{ka:tbl_zn} приведены параметры, необходимые для настройки функционального блока <<ЗН>>.

\vspace{2mm}
\begin{figure}[H]
\centering
\includegraphics[width=0.9\textwidth,height=0.9\textheight,keepaspectratio]{\fbpath/213. КА ТН/_latex/img/zn.pdf}
\captionof{figure}{Функциональная схема <<ЗН>>}\label{ka:img_zn}
\end{figure}

\small
\begin{longtable}{|>{\centering\arraybackslash}m{5.3cm}|>{\centering\arraybackslash}m{3.3cm}|>{\centering\arraybackslash}m{4.2cm}|>{\centering\arraybackslash}m{1.8cm}|>{\centering\arraybackslash}m{1cm}|}
\caption{Параметры для настройки функции <<ЗН>>\hfill\vspace{-0.5\baselineskip}}\label{ka:tbl_zn}\\ 
\hline
\rowcolor{gray!30}
Параметр (Параметр на ИЧМ) & Условное обозначение на схеме & Значение/ Диапазон & Единица измерения & Шаг \\ 
\hline
\endfirsthead
\caption*{\hspace{3pt}\emph{Продолжение таблицы \ref{ka:tbl_zn}\hfill\vspace{-0.5\baselineskip}}} \\ % сделано по ГОСТ 2.105 п.6.8.7
\hline
\rowcolor{gray!30}
Параметр (Параметр на ИЧМ) & Условное обозначение на схеме & Значение/ Диапазон & Единица измерения & Шаг \\ 
%\hline
\endhead
%\hline
\endfoot
%\hline
\endlastfoot
%==+t1*ES1|VT_SwitchDevice
\centering Ввод функции в работу (Ввод\_функции) & \centering SGF1 & \centering 0 = Не предусмотрено\\1 = Предусмотрено & \centering -- & \centering \arraybackslash -- \\
\hline
\centering Режим отключения ЗН (Режим\_откл) & \centering SGF2 & \centering 0 = Длительный\\1 = Импульсный & \centering -- & \centering \arraybackslash -- \\
\hline
\centering Режим включения ЗН (Режим\_вкл) & \centering SGF3 & \centering 0 = Длительный\\1 = Импульсный & \centering -- & \centering \arraybackslash -- \\
\hline
\centering Длительность формирования сигнала о неисправном положении ЗН (Tнеиспр\_ЗН) & \centering T1 & \centering 10 ... 100000 & \centering мс & \centering \arraybackslash 10 \\
\hline
\centering Длительность формирования команды «Отключить ЗН (реле)» (Тимп\_откл) & \centering T2 & \centering 100 ... 100000 & \centering мс & \centering \arraybackslash 10 \\
\hline
\centering Длительность формирования команды «Включить ЗН (реле)» (Тимп\_вкл) & \centering T3 & \centering 100 ... 100000 & \centering мс & \centering \arraybackslash 10 \\
\hline
%===t1
\end{longtable}
\normalsize

\end{enumerate}

%%%%%%%%%%%%%%%%%%%%%%%%%%%%%%%%%%%%%%%%%%%%%%%%%%%%%%%%%%% ЗН %%%%%%%%%%%%%%%%%%%%%%%%%%%%%%%%%%%%%%%%%%%%%%%%%%%%%%%%%%

\item Выкатной элемент (ВЭ)

\begin{enumerate}[label=\arabic{section}.\arabic{subsection}.\arabic{enumi}.\arabic*, labelsep=4pt, leftmargin=0em, itemindent=65pt, parsep=0pt]

\item
При наличии входного сигнала <<УВЭ: Выкатить>> будет формироваться команда выкатить выкатной элемент <<ВЭ: Выкатить (реле)>> на внутреннее выходное реле устройства с целью коммутации выходной цепи выкатного элемента. С помощью уставки SGF2 <<Режим\_выкат>> задается режим работы команды выкатить ВЭ (длительный/импульсный). В импульсном режиме, выдаваемый импульс, длится на протяжении времени, задаваемого уставкой Т2. При наличии контрольного положения ВЭ произойдет сброс команды выкатить ВЭ как в импульсом, так и в длительном режимах. Сброс команды выкатить также осуществляется с помощью команды <<Сброс>>.
\item
При наличии входного сигнала <<УВЭ: Вкатить>> будет формироваться команда вкатить выкатной элемент <<ВЭ: Вкатить (реле)>> на внутреннее выходное реле устройства с целью коммутации выходной цепи выкатного элемента. С помощью уставки SGF3 <<Режим\_вкат>> задается режим работы команды вкатить ВЭ (длительный/импульсный). В импульсном режиме, выдаваемый импульс, длится на протяжении времени, задаваемого уставкой Т3. При наличии рабочего положения ВЭ произойдет сброс команды вкатить ВЭ как в импульсом, так и в длительном режимах. Сброс команды вкатить также осуществляется с помощью команды <<Сброс>>.

\item
При выводе функции из работы (отсутствие выходного сигнала <<ВЭ: Ввод>>) происходит сброс команд управления выкатным элементом.

\item
В логике функции реализовано формирование выходных сигналов от блок-контактов тележки ВЭ без выдержки времени:
\begin{itemize}
    \item сигнал промежуточного положения тележки ВЭ <<ВЭ: Промеж. положение>> формируется, когда происходит перемещения тележки из рабочего положения в ремонтное (контрольное) или обратно. При этом действии происходит одновременное наличие (отсутствие) обоих сигналов от блок-контактов тележки ВЭ <<Контр.полож. ВЭ>>, <<Рабоч.полож. ВЭ>>;
    \item когда ВЭ находится в контрольном положении или полностью выкачен в ремонтное положение формируется сигнал <<ВЭ: Выкачен>>;
    \item сигнал <<ВЭ: Вкачен>> формируется, когда тележка ВЭ находится в рабочем положении.
\end{itemize}
\item
Неисправное положение ВЭ <<ВЭ: Промеж. положение>> формируется по истечении регулируемой выдержки времени задаваемой уставкой таймера Т1 и наличии подключенного в разъем шлейфа управления.
\item
Уставкой SGF4 <<Контр\_рем\_полож>> контролируется ремонтное положение ВЭ. Если шлейф управления извлечен из разъема, то на выходе формируется сигнал, что ВЭ находится в ремонтном положении <<ВЭ: Ремонт. положение>>.

\item
Ввод функции <<ВЭ>> в работу на этапе параметрирования устройства осуществляется программным переключателем SGF1 <<Ввод\_функции>>. Введенное состояние функции соответствует активному состоянию сигнала <<ВЭ: Ввод>>.

\item
Алгоритм работы функции <<ВЭ>> выполнен в соответствии с рисунком \ref{ka:img_ve}. В таблице \ref{ka:tbl_ve} приведены параметры, необходимые для настройки функционального блока <<ВЭ>>.

\vspace{2mm}
\begin{figure}[H]
\centering
\includegraphics[width=0.9\textwidth,height=0.9\textheight,keepaspectratio]{\fbpath/213. КА ТН/_latex/img/ve.pdf}
\captionof{figure}{Функциональная схема <<ВЭ>>}\label{ka:img_ve}
\end{figure}

\small
\begin{longtable}{|>{\centering\arraybackslash}m{5.3cm}|>{\centering\arraybackslash}m{3.3cm}|>{\centering\arraybackslash}m{4.2cm}|>{\centering\arraybackslash}m{1.8cm}|>{\centering\arraybackslash}m{1cm}|}
\caption{Параметры для настройки функции <<ВЭ>>\hfill\vspace{-0.5\baselineskip}}\label{ka:tbl_ve}\\ 
\hline
\rowcolor{gray!30}
Параметр (Параметр на ИЧМ) & Условное обозначение на схеме & Значение/ Диапазон & Единица измерения & Шаг \\ 
\hline
\endfirsthead
\caption*{\hspace{3pt}\emph{Продолжение таблицы \ref{ka:tbl_ve}\hfill\vspace{-0.5\baselineskip}}} \\ % сделано по ГОСТ 2.105 п.6.8.7
\hline
\rowcolor{gray!30}
Параметр (Параметр на ИЧМ) & Условное обозначение на схеме & Значение/ Диапазон & Единица измерения & Шаг \\ 
%\hline
\endhead
%\hline
\endfoot
%\hline
\endlastfoot
%==+t1*WE1|VT_SwitchDevice
\centering Ввод функции в работу (Ввод\_функции) & \centering SGF1 & \centering 0 = Не предусмотрено\\1 = Предусмотрено & \centering -- & \centering \arraybackslash -- \\
\hline
\centering Режим работы (выкатить) ВЭ (Режим\_выкат) & \centering SGF2 & \centering 0 = Длительный\\1 = Импульсный & \centering -- & \centering \arraybackslash -- \\
\hline
\centering Режим работы (вкатить) ВЭ (Режим\_вкат) & \centering SGF3 & \centering 0 = Длительный\\1 = Импульсный & \centering -- & \centering \arraybackslash -- \\
\hline
\centering Контроль ремонтного положения (Контр\_рем\_полож) & \centering SGF4 & \centering 0 = Не предусмотрено\\1 = Предусмотрено & \centering -- & \centering \arraybackslash -- \\
\hline
\centering Длительность формирования сигнала о неисправном положении ВЭ (Tнеиспр\_ВЭ) & \centering T1 & \centering 10 ... 100000 & \centering мс & \centering \arraybackslash 10 \\
\hline
\centering Длительность формирования команды «Выкатить ВЭ (реле)» (Тимп\_выкат) & \centering T2 & \centering 100 ... 100000 & \centering мс & \centering \arraybackslash 10 \\
\hline
\centering Длительность формирования команды «Вкатить ВЭ (реле)» (Тимп\_вкат) & \centering T3 & \centering 100 ... 100000 & \centering мс & \centering \arraybackslash 10 \\
\hline
%===t1
\end{longtable}
\normalsize

\end{enumerate}

\end{enumerate}
